% The Persistence of Structure
% From Shape Space to Recursive Persistence:
% Julian Barbour, Entaxy, and the RSVP Framework
% LuaLaTeX — TeX Gyre Pagella + unicode-math

\documentclass[12pt, a4paper]{article}

%% ── Font & encoding ────────────────────────────────────────────────────────
\usepackage{fontspec}
\usepackage[math-style=ISO]{unicode-math}

\setmainfont{TeX Gyre Pagella}
\setsansfont{TeX Gyre Heros}[Scale=MatchLowercase]
\setmathfont{Latin Modern Math}[Scale=MatchLowercase]

%% ── Micro-typography & spacing ─────────────────────────────────────────────
\usepackage{microtype}
\usepackage{setspace}
\setstretch{1.18}

\usepackage[skip=8pt plus 2pt, indent=0pt]{parskip}

%% ── Page geometry ──────────────────────────────────────────────────────────
\usepackage[
  a4paper,
  top=38mm,
  bottom=38mm,
  left=36mm,
  right=36mm,
  headheight=14pt
]{geometry}

%% ── Headers / footers ──────────────────────────────────────────────────────
\usepackage{fancyhdr}
\pagestyle{fancy}
\fancyhf{}
\renewcommand{\headrulewidth}{0.3pt}
\fancyhead[L]{\small\textit{The Persistence of Structure}}
\fancyhead[R]{\small\textit{Flyxion}}
\fancyfoot[C]{\small\thepage}

%% ── Section styling ────────────────────────────────────────────────────────
\usepackage{titlesec}
\titleformat{\section}
  {\normalfont\large\bfseries}
  {\Roman{section}.}{0.8em}{}[\vspace{-2pt}\rule{\linewidth}{0.3pt}\vspace{2pt}]
\titleformat{\subsection}
  {\normalfont\normalsize\itshape}
  {}{0em}{}
\titlespacing*{\section}{0pt}{18pt}{8pt}
\titlespacing*{\subsection}{0pt}{10pt}{4pt}

%% ── Hyperlinks ─────────────────────────────────────────────────────────────
\usepackage[
  colorlinks=true,
  linkcolor=black,
  citecolor=black,
  urlcolor=black
]{hyperref}

%% ── Miscellaneous ──────────────────────────────────────────────────────────
\usepackage{xcolor}
\usepackage{booktabs}
\usepackage{enumitem}
\setlist[itemize]{leftmargin=1.4em, itemsep=2pt, topsep=4pt}

%% ── TikZ ───────────────────────────────────────────────────────────────────
\usepackage{tikz}
\usetikzlibrary{arrows.meta, positioning, calc, decorations.pathmorphing,
  shapes.geometric, fit}
\usepackage{float}

%% ── Custom operators / symbols ─────────────────────────────────────────────
% Admissibility field, projection, recursive consistency, manifold
\newcommand{\Afield}{\mathcal{A}}
\newcommand{\proj}{\hat{\pi}}
\newcommand{\Rcomp}{\mathcal{R} \circ \mathcal{R}}
\newcommand{\Madm}{\mathbb{M}_{a}}
\newcommand{\Xtraj}{X}
\newcommand{\admlog}{\mathcal{L}_{\mathrm{adm}}}

%% ══════════════════════════════════════════════════════════════════════════
\begin{document}

%% ── Title page ─────────────────────────────────────────────────────────────
\begin{titlepage}
\thispagestyle{empty}
\vspace*{14mm}

{\centering
  {\fontsize{17}{21}\selectfont\bfseries The Persistence of Structure}\\[8pt]
  {\normalsize\itshape From Shape Space to Recursive Persistence:\\[3pt]
  Julian Barbour, Entaxy, and the RSVP Framework}\\[12mm]

  {\normalsize Flyxion}\\[4pt]
  {\small Independent Researcher}\\[3pt]
  {\small\color{gray} \today}\\[10mm]
\par}

\begin{abstract}
\noindent
Julian Barbour's relational cosmology, developed most fully in \textit{The Janus Point},
represents a profound ontological inversion: time, scale, and reference frames dissolve
into evolving relational configurations, leaving complexity growth as the universe's
dominant arrow. Yet Barbour's framework remains transitional. Its dynamical grammar
is still inherited from Newtonian mechanics, complexity is tracked but not
explained operationally, and the Janus Point functions as a hidden global anchor
despite the theory's Machian intentions. This essay argues that the Relativistic
Scalar-Vector Plenum (RSVP) extends Barbour's relational reduction into a
fully recursive consistency framework. In the RSVP picture, the physically
meaningful quantity is not entropy, complexity, or entaxy, but
\emph{recursively maintainable coherence}: only configurations that preserve
admissibility across scales persist long enough to constitute stable structure.
Barbour's Janus Point becomes a local admissibility extremum rather than a
universal origin; his shape space becomes a special case of a constraint-selective
projection manifold; and his records are reinterpreted as active operational
constraints required for future inferential continuity. The essay concludes by
extending the framework toward cognition, where the same recursive admissibility
operators that govern cosmological persistence reappear as the structural
conditions for memory, measurement, and predictive continuity.
\end{abstract}

\end{titlepage}

\newpage
\setcounter{page}{1}

%% ══════════════════════════════════════════════════════════════════════════
\section{Introduction: The Ontological Inversion}

The history of physics can be read as a sequence of ontological compressions.
Classical mechanics treated reality as matter evolving inside a pre-existing
spatial container governed by externally imposed laws. Relativity destabilized
the fixed background, merging space and time into dynamical geometry, yet
retained the assumption that lawful structure exists independently of the
systems it governs. Modern cosmology inherited both frameworks while adding
statistical thermodynamics, producing the contemporary picture of an
expanding universe in which entropy increase provides the primary directional
arrow.

Julian Barbour's relational cosmology represents one of the most sustained
attempts to overturn this inheritance. In \textit{The Janus Point} (2020),
Barbour argues that neither absolute space, absolute scale, nor conventional
time are fundamental. Instead, the universe should be described entirely in
terms of evolving relational configurations. The universe does not expand
through an external metric backdrop; its internal relational geometry changes
\emph{shape}. Time itself becomes reconstructible from ordered relational
differences rather than existing as an independent parameter.

This constitutes a genuine ontological inversion. Reality ceases to be a
collection of objects embedded within spacetime and becomes instead a
trajectory through relational configuration space. The Machian commitment is
total: change is primary; everything else is abstracted from it.

Yet Barbour's framework remains transitional. While it successfully removes
external coordinate structures, it still presupposes a globally coherent
dynamical grammar governing the evolution of shape space. The Janus Point
functions as a privileged geometric pivot around which temporal asymmetry
organizes. Complexity increases statistically as the system evolves away from
this pivot, but the framework does not fully account for why coherent
structure \emph{persists} rather than collapsing into rigidity or
dissipating into noise.

The Relativistic Scalar-Vector Plenum (RSVP) extends the relational turn
beyond geometry into recursive admissibility. Rather than treating physical
law as a pre-existing grammar operating on configurations, RSVP interprets
persistence itself as the result of constraint-selection across admissibility
landscapes. Stable structures survive not merely because they are
statistically favored, but because they maintain recursive consistency across
scales.

The transition from Barbour to RSVP therefore marks a deeper inversion:

The transition moves from relational geometry to operational persistence,
from lawful evolution to admissibility management, and from static ontology
to recursive consistency architecture.

The Janus Point becomes not the origin of order, but a special case within a
broader landscape of recursive stabilization.

%% ══════════════════════════════════════════════════════════════════════════
\section{The Geometry of the Minimal State}

\subsection{Barbour's Global Anchor}

The Janus Point emerges from the Newtonian \(N\)-body problem under conditions
of non-negative total energy. In such systems, the global scale of the
particle distribution decreases from infinity to a finite minimum before
increasing again. Near this minimal configuration, particle distributions
become maximally uniform — the most homogeneous state the system can occupy.
Moving away from the Janus Point in either temporal direction produces
increasing clustering and relational differentiation.

Barbour formulates this through the concept of \emph{complexity}: a
scale-invariant quantity measuring the degree to which particles cluster
rather than distribute uniformly. At the Janus Point, complexity is minimal.
As the system evolves outward, complexity grows monotonically in both
temporal directions, giving rise to two bidirectional arrows of time that
point away from the central configuration.

The elegance of this picture lies precisely in its avoidance of imposed
initial conditions. Temporal asymmetry does not need to be stipulated; it
emerges from the gravitational dynamics of relational structure.

\subsection{The Problem of the Hidden Anchor}

Nevertheless, the Janus Point introduces a structural tension that becomes
visible from within an admissibility framework.

Although Barbour explicitly rejects absolute time, the Janus Point functions
as a globally privileged geometric state. It implicitly establishes a
universal ordering principle — a hidden absolute anchor — around which the
entire cosmological trajectory is coordinated. Even if no observer can access
both halves of the resulting bi-universe simultaneously, the theory still
relies upon a globally distinguished minimum configuration as its
organizational center.

The theory therefore smuggles a form of cosmic meta-time ordering back into
the geometry. The Janus Point is defined precisely as the moment when global
shape has reached its most uniform state, which means the theory already
presupposes that "global state" is a coherent and uniquely identifiable
object.

That presupposition is not trivially available in a fully relational framework.

\subsection{Decentralized Admissibility}

RSVP critiques this requirement for a singular global minimum.

In recursive admissibility frameworks, there is no necessity for one
privileged universal pivot. Instead, local regions of high-symmetry
reintegration continuously emerge within the plenum as part of ongoing
dynamical regulation. Uniformity is not a singular historical event but a
recurring accessibility condition — something the system negotiates locally
rather than instantiates once.

The cosmological question shifts accordingly. Rather than asking why the
universe began in a low-entropy state, RSVP asks:

\begin{quote}
\itshape
What conditions permit recursively maintainable structure to emerge and
stabilize within an admissibility landscape?
\end{quote}

This dissolves the need for a universal temporal pivot while preserving the
emergence of directional persistence. Janus-like states become local
admissibility extrema — recurring features of the landscape rather than
cosmological origins.

\begin{figure}[H]
\centering
\begin{tikzpicture}[scale=0.95]
  % Axis
  \draw[thick] (-5.2,0) -- (5.2,0);

  % Janus point
  \filldraw (0,0) circle (2.5pt);
  \node[below=4pt] at (0,0) {\small Janus Point};

  % Bidirectional time arrows
  \draw[-{Latex[length=3mm]}, thick] (0,0) -- (-4.8,0);
  \draw[-{Latex[length=3mm]}, thick] (0,0) -- (4.8,0);
  \node[below=4pt] at (-3.5,0) {\small \textit{t} $\to -\infty$};
  \node[below=4pt] at (3.5,0) {\small \textit{t} $\to +\infty$};

  % Complexity growth curves
  \draw[thick] (0,0) .. controls (1,0.9) and (2.5,1.6) .. (4.8,2.0);
  \draw[thick] (0,0) .. controls (-1,0.9) and (-2.5,1.6) .. (-4.8,2.0);
  \node[right] at (4.8,2.0) {\small complexity};

  % RSVP local basins (dashed circles at distributed positions)
  \foreach \x/\y in {-3.8/1.7, -1.6/0.8, 1.4/0.7, 3.5/1.6} {
    \draw[dashed, gray] (\x,\y) circle (0.45);
    \filldraw[gray] (\x,\y) circle (1.8pt);
  }
  \node[above=2pt] at (-3.8,2.15) {\small\itshape local basin};
  \node[above=2pt] at (3.5,2.05) {\small\itshape local basin};
\end{tikzpicture}
\caption{Barbour's bidirectional Janus geometry (solid) contrasted with
RSVP distributed local admissibility extrema (dashed). RSVP requires no
privileged global pivot; admissibility basins recur across the plenum.}
\end{figure}

%% ══════════════════════════════════════════════════════════════════════════
\section{Complexity and the Arrow of Time}

\subsection{The Boltzmann Box and Its Limits}

Barbour's critique of entropy centers on confinement. Classical
thermodynamics emerged from systems enclosed within boundaries — steam
engines, gases in containers, mechanically isolated volumes. Under such
conditions, systems evolve toward equilibrium distributions, and entropy
becomes associated with increasing uniformity.

Barbour argues this intuition fails entirely for unconfined gravitational
systems. Gravity naturally amplifies clustering rather than suppressing it.
A gas in a box spreads; a gravitational system collapses into structure.
Applying box-thermodynamics to cosmology therefore produces what Barbour
regards as a systematic distortion: the universe's increasing complexity gets
misread as entropy increase rather than recognized as relational
differentiation.

To formalize this inversion, Barbour introduces \emph{entaxy} — a measure
of the admissible volume of relational configuration space corresponding to
a given level of structural organization. As the universe moves away from
the Janus Point, entaxy decreases while complexity increases. Structure is
not disorder; it is organized relational specificity.

This is a strong and largely correct criticism of naive thermodynamic
cosmology. Gravitational entropy has always been conceptually anomalous
precisely because gravity produces structure rather than diffusion. Barbour
geometrizes this anomaly rather than dismissing it.

\subsection{The Limits of Entaxy}

Yet entaxy remains fundamentally descriptive.

It explains why clustering emerges statistically in unconfined gravitational
systems. It does not explain why coherent structure persists dynamically —
why some configurations remain recursively stable across scales while others
collapse, decohere, or dissipate. Entaxy tracks the growth of structural
differentiation, but it provides no mechanism by which the system avoids
two failure modes: terminal rigidity (over-concentrated, causally isolated
clusters) and diffusive dissolution (structures that differentiate but
cannot maintain internal coherence).

This is the central gap between Barbour's framework and an admissibility
theory.

\subsection{The Lamphron–Lamphrodyne Duality}

RSVP addresses this gap through the lamphron–lamphrodyne duality — two
operationally defined pressures governing the dynamics of structural
persistence.

\textbf{Lamphron} designates the differentiation pressure: the tendency for
systems to generate localized structure, constrain relational identity, form
stable memory configurations, and specialize toward predictive specificity.
Operationally, lamphron corresponds to local trajectory specialization that
increases admissibility resolution — the system narrows its accessible future
states by committing to particular relational configurations.

\textbf{Lamphrodyne} designates the reintegration pressure: the
redistributive and smoothing dynamics that prevent terminal rigidity.
Operationally, lamphrodyne corresponds to the entropy-redistribution
dynamics required when a configuration becomes recursively overconstrained
— when its internal coherence can no longer be maintained without external
admissibility exchange. A configuration under lamphrodyne pressure must
redistribute information and energy across scales to remain dynamically
legible.

Together, these pressures produce what might be called a \emph{breathing}
cosmology. Unlike Barbour's one-way trajectory from uniformity to
clustering, RSVP describes an ongoing cycle in which structure is
continuously vetted against its capacity to remain recursively admissible.

\subsection{Persistence as the Physically Meaningful Quantity}

This leads to the central claim distinguishing RSVP from both classical
thermodynamics and Barbour's entaxy framework:

\begin{quote}
\itshape
Only recursively admissible configurations persist.
\end{quote}

The arrow of time therefore ceases to be merely the direction of increasing
clustering, complexity, or entaxy decrease. It becomes the direction in
which predictive continuity remains constructible — the direction in which
the lamphron–lamphrodyne balance permits ongoing structural stabilization
rather than terminal collapse or dissolution.

This separates three distinct pictures of cosmological directionality:

Standard thermodynamics holds that systems evolve toward equilibrium;
Barbour's framework holds that systems evolve toward clustering complexity;
RSVP holds that systems evolve toward recursively maintainable admissibility
regimes. These are not equivalent descriptions. A crystal is highly ordered; a star
is dynamically structured; a hurricane is locally coherent; a biosphere is
recursively stabilized. Only the last constitutes persistence in the RSVP
sense.

%% ══════════════════════════════════════════════════════════════════════════
\section{The Emergence of Reference: Projection and Constraint}

\subsection{Kepler Pairs as Emergent Clocks}

One of Barbour's most important insights concerns the genesis of
measurement. As gravitational systems evolve away from the Janus Point,
particles form stable Kepler pairs — gravitationally bound subsystems that
function simultaneously as rods, clocks, and compasses. Measurement
therefore emerges from dynamical stabilization rather than being imposed
from outside.

This is a profound observation. Clocks are not coordinate devices
superimposed on physics. They are dynamically stabilized relational
attractors — structures that persist long enough to serve as reference
baselines for other structures.

\subsection{The Projection Formalism}

RSVP formalizes this through the projection operator:

\[
  \proj : \Xtraj \;\longrightarrow\; \Madm
\]

Here \(\Xtraj\) represents the underlying high-dimensional trajectory space
of all dynamically possible paths, while \(\Madm\) represents the compressed
admissibility manifold — the set of trajectories that remain operationally
accessible given the system's current constraint structure.

Barbour's shape space can be interpreted as an early relational form of such
a projection geometry. His central move — quotienting out translations,
rotations, and absolute scale to retain only relational configurations — is
mathematically close to the RSVP compression operator. Both frameworks
reduce redundant coordinates and preserve only dynamically meaningful
relational structure.

\subsection{Projection as Cosmological Grammar}

However, RSVP introduces a claim that Barbour's framework does not make
explicit: projections are dynamically consequential. Compression changes
future admissibility conditions.

Consider observational cosmology itself. Once the universe is interpreted
through the grammar of metric expansion, subsequent theoretical trajectories
become trapped within expansion-compatible explanatory manifolds. Problems
such as dark matter, dark energy, inflation, and the horizon problem then
appear as downstream stabilization artifacts of the original projection
choice — attempts to preserve coherence within the constraints set by the
compressed representation.

This is not a criticism of metric expansion as such. It is an illustration
of a general principle: the choice of projection manifold determines which
futures remain conceptually accessible. A measuring device, a coordinate
system, and a theoretical framework all function as projection operators
that constrain subsequent admissibility.

Measurement does not passively record reality. It participates in its
stabilization. Reference frames are not descriptive tools; they are active
operators of recursive persistence, trapping future trajectories into
specific admissibility namespaces.

\begin{figure}[H]
\centering
\begin{tikzpicture}[
    node distance=2.6cm,
    every node/.style={align=center},
    box/.style={draw, rectangle, minimum width=3cm, minimum height=1cm,
                font=\small}
]
  \node[box] (X) {\(\Xtraj\)\\[2pt]Trajectory Space};
  \node[box, right=of X] (M) {\(\Madm\)\\[2pt]Admissibility Manifold};
  \node[box, dashed, below=2cm of M] (F)
        {Expansion-Compatible\\Namespace};

  \draw[-{Latex[length=3mm]}, thick] (X) -- node[above]{\(\proj\)} (M);
  \draw[-{Latex[length=3mm]}, thick] (M) -- (F);
  \draw[dotted, thick] (X) to[bend right=30]
        node[below, font=\small\itshape]{inaccessible futures} (F);
\end{tikzpicture}
\caption{The projection operator \(\proj\) compresses total trajectory space
\(\Xtraj\) into the admissibility manifold \(\Madm\), constraining which
future explanatory structures remain accessible. The dashed box illustrates
how a particular projection grammar (here, metric expansion) traps subsequent
theoretical development within a specific namespace.}
\end{figure}

%% ══════════════════════════════════════════════════════════════════════════
\section{The Localist Turn: Event-Level Admissibility}

\subsection{Against Global Reconstruction}

One of the deepest assumptions shared by both conventional signal processing
and conventional cosmology is the assumption of global reconstruction.
Raw events are aggregated into large representational buffers and only then
transformed into coherent structure through global operations. In signal
processing this appears through Fourier analysis and spectral decomposition.
In cosmology it appears through globally coordinated metric descriptions,
universal state functions, and reconstruction of cosmic history through
large-scale geometric inference.

This methodology implicitly assumes that coherence is something recovered
\emph{after} aggregation.

RSVP challenges this assumption at the architectural level.

In recursively admissible systems, coherence already exists at the event
layer itself. Stable recurrence patterns do not need to be globally solved
because admissible timing relations are locally preserved through recursive
constraint filtering. The system does not reconstruct coherent structure
through a global transform; it recognizes structure through persistent local
parity, interval stability, and admissible recurrence.

The distinction is not merely computational. It is ontological.

A representation-first architecture assumes that reality becomes meaningful
only after sufficient global aggregation. A persistence-first architecture
assumes that meaningful structure emerges directly through local recursive
stabilization.

\subsection{Shape Space and Local Persistence}

This clarifies the central divergence between Barbour's relational cosmology
and the RSVP framework.

Barbour's shape dynamics still presupposes a globally coherent trajectory
through relational configuration space. Complexity is defined through the
global organization of the system as a whole, and the Janus Point functions
as a distinguished geometric extremum coordinating the emergence of temporal
directionality.

RSVP instead treats large-scale structure as the emergent consequence of
distributed local admissibility relations. Geometry is not solved globally
and then imposed downward. It emerges upward from recursively stabilized
constraint structures operating directly at the event level.

This transforms the role of cosmological coordination entirely. Reference
frames become residual stabilization artifacts rather than pre-existing
coordinate systems. Cosmological regularity becomes the consequence of
distributed admissibility filtering rather than globally synchronized
evolution. Memory becomes a recursively accessible trajectory preserved by
operational continuity rather than a passive storage of representations.

The universe therefore behaves less like a globally integrated equation and
more like a locality-preserving recursive filter.

\begin{figure}[H]
\centering
\begin{tikzpicture}[scale=0.7]
  % Background grid
  \foreach \x in {0,...,7} {
    \foreach \y in {0,...,7} {
      \draw[gray!40] (\x,\y) rectangle ++(1,1);
    }
  }
  % Central event
  \filldraw[black] (3.5,3.5) circle (5pt);
  \node[above=6pt, font=\small] at (3.5,3.5) {event};
  % Propagation ripples
  \draw[dashed, gray!70] (3.5,3.5) circle (1.3);
  \draw[dashed, gray!50] (3.5,3.5) circle (2.2);
  \draw[dashed, gray!30] (3.5,3.5) circle (3.1);
  % Local stabilization arrows
  \draw[-{Latex[length=2mm]}, thick] (3.5,3.5) -- (4.5,4.3);
  \draw[-{Latex[length=2mm]}, thick] (3.5,3.5) -- (2.4,4.2);
  \draw[-{Latex[length=2mm]}, thick] (3.5,3.5) -- (4.4,2.6);
  \draw[-{Latex[length=2mm]}, thick] (3.5,3.5) -- (2.5,2.7);
  % Labels
  \node[font=\small\itshape] at (6.5,7.2) {admissibility};
  \node[font=\small\itshape] at (6.5,6.7) {propagation};
\end{tikzpicture}
\caption{Local event-level admissibility propagation. Coherence emerges
through recursive neighbourhood stabilization rather than global
reconstruction. Dashed rings indicate successive admissibility horizons;
arrows indicate locally constrained trajectory directions.}
\end{figure}

\subsection{Persistence as the Fundamental Invariant}

The resulting shift is profound. Global descriptions do not disappear because
they are false; they become secondary accounts of structures that are already
locally stabilized.

The primary invariant is no longer global geometry but recursive persistence
itself.

Structures survive because they remain admissible across scales. Local
constraint systems continuously negotiate stabilization without requiring a
single centralized reconstruction of total state. The cosmos persists not
because it is coordinated by a universal supervisory geometry, but because
distributed recursive consistency-management prevents admissible structures
from vanishing faster than they can reintegrate.

This provides a very different image of physical law. Laws cease to appear
as opaque global commands imposed upon matter from outside. They become
emergent regularities generated by locality-preserving recursive persistence.

The universe becomes legible precisely because persistence remains locally
repairable.

%% ══════════════════════════════════════════════════════════════════════════
\section{Records as Operational Constraints}

\subsection{Barbour's Geometric Records}

Barbour frequently treats records as fossilized relational traces embedded
within present configurations. The past survives through stable geometric
correlations preserved in current arrangements of matter. Memory is a kind
of structural sediment — the present configuration carries the imprint of
prior shapes.

This is geometrically elegant and philosophically significant. It dissolves
the traditional asymmetry between past and future by locating temporal
evidence entirely within present structure. There is no need for a
metaphysical "passage of time"; only configurations that encode differential
information about prior states.

\subsection{The Admissibility Log}

RSVP radicalizes this insight in a direction Barbour does not pursue.

Records are not merely traces of prior states. They are operational
constraints required for future continuity. The Admissibility Log,
\(\admlog\), formalizes this principle: a record exists not because it
describes a past event, but because maintaining it preserves recursive
consistency across future trajectories.

The distinction matters enormously.

In Barbour's account, a record is a geometric artifact — something left
behind by prior dynamical evolution, epistemically available to present
observers. In the RSVP account, a record is a \emph{functional constraint}
— something the system must maintain in order to remain recursively
admissible. A record that ceases to be operationally compatible with the
system's future trajectory is not merely lost information; it is a structural
failure that changes which futures remain accessible.

Memory, on this account, is not retrospective. It is infrastructural.

\subsection{Time as Directional Record-Utility}

This reframes the nature of time itself.

Time is not fundamentally the accumulation of complexity. Nor is it the
ordering of successive relational shapes. Time, in the RSVP picture,
becomes the directional persistence of usable records — the direction in
which prior constraints remain operationally functional for maintaining
future admissibility.

A system experiences temporal continuity insofar as its admissibility log
remains coherently extensible. When prior constraints can no longer be
integrated into future trajectories — when the records lose operational
compatibility — the system undergoes a form of structural rupture that
constitutes, in a precise sense, temporal discontinuity.

Causality follows from this picture. Causes are not merely prior events.
They are persistent admissibility constraints governing which futures remain
accessible. The arrow of time is therefore the direction of
record-constructibility: the direction in which the admissibility log
supports rather than blocks future trajectory space.

%% ══════════════════════════════════════════════════════════════════════════
\section{Cognition as Recursive Admissibility}

\subsection{From Cosmology to Cognitive Architecture}

The transition from cosmological admissibility to cognitive architecture is
not a change of subject. It is the recognition that the same organizational
operators recur across scales because persistence itself obeys shared
accessibility constraints.

Consider what a cognitive system must do to remain operationally functional.
It must maintain an admissibility log — a record structure that preserves
inferential continuity across successive states. It must project
high-dimensional environmental trajectories onto compressed manifolds that
remain actionable. It must balance lamphron pressure (the drive toward
increasingly specific predictive representations) against lamphrodyne
pressure (the need to remain flexible enough to revise constraints without
structural collapse).

A mind, on this account, is not a container of representations. It is a
recursively stabilized admissibility manifold — a system that persists by
continuously negotiating the compression-reintegration balance.

\subsection{Memory, Measurement, and Predictive Continuity}

Barbour's Kepler pairs and RSVP's projection operators converge here in
a striking way. Both treat measurement as emergent stabilization rather
than external imposition. Both recognize that clocks, rods, and reference
frames are persistent relational attractors rather than coordinate
conventions.

In cognitive terms: perceptual systems are not passive recorders. They are
projection operators that collapse environmental trajectory space into
admissible action manifolds. Memory is not archival storage; it is the
operational constraint structure that keeps future inference admissible.
Attention is lamphron pressure — the selective specialization that increases
predictive resolution at the cost of narrowing accessible futures.

The pathological failure modes are symmetric. Over-specialization (excessive
lamphron pressure) produces rigid, non-revisable predictive models —
cognitive brittleness. Over-generalization (excessive lamphrodyne pressure)
produces diffuse, incoherent representations incapable of maintaining
predictive continuity — cognitive dissolution.

The healthy cognitive system, like the healthy cosmological structure, is one
that maintains the lamphron–lamphrodyne balance across scales — recursively
admissible at every level of description.

%% ══════════════════════════════════════════════════════════════════════════
\section{Against the Janus Reversal}

\subsection{The Problem of Bidirectional Time}

One of the most philosophically striking claims in Barbour's framework is
that temporal arrows emerge bidirectionally from the Janus Point. As the
system evolves away from the point of minimal complexity, observers on both
sides experience what appears locally to be a forward-moving arrow of time.
The resulting picture is not a multiverse but what Barbour calls a
\emph{bi-universe}: two temporally opposed branches emerging from a common
state of maximal uniformity.

The elegance of this construction is undeniable. It preserves the time
symmetry of the underlying equations while simultaneously accounting for the
experienced asymmetry of temporal direction. Yet the construction introduces
a deep conceptual difficulty.

The Janus reversal presupposes that the universe possesses a globally
coherent trajectory capable of being meaningfully partitioned into two
temporal halves. The universe must therefore already exist as a unified
meta-history before either branch can inherit its local arrow.

In effect, the framework removes absolute time only to reintroduce a hidden
global chronology through geometric structure itself. The Janus Point is
defined as the moment when global shape reaches maximal uniformity — but
this already presupposes that ``global state'' is a uniquely identifiable and
coherent object. That presupposition is not trivially available in a fully
relational framework.

RSVP rejects this requirement. Time does not reverse at a Janus point
because recursive admissibility does not require a globally mirrored
trajectory. Temporal direction emerges locally from the preservation of
predictive continuity within recursively stabilized structures. The arrow of
time is therefore not generated by motion away from a universal minimum, but
by the directional extension of admissibility-preserving records.

The distinction is precise:
\begin{align*}
\text{Barbour:}&\quad \text{Minimal Complexity} \;\longrightarrow\; \text{Bidirectional Time} \\
\text{RSVP:}&\quad \text{Recursive Persistence} \;\longrightarrow\; \text{Directional Continuity}
\end{align*}

Temporal asymmetry does not arise because the universe expands away from a
special point. It arises because recursively admissible structures can only
maintain coherent predictive extension in one operational direction at a
time.

\subsection{Persistence Without Reversal}

The RSVP framework therefore removes the need for a cosmological mirror
branch entirely. There is no requirement for time to face both ways from a
privileged geometric pivot because admissibility itself already generates
directionality.

A record either remains operationally extensible or it does not. A
constraint either preserves predictive continuity or it fails. The arrow of
time is simply the direction in which admissibility logs remain recursively
constructible.

This removes one of the most problematic implications of Janus cosmology:
that the universe requires a globally coordinated temporal bifurcation to
produce local asymmetry. RSVP instead treats temporal asymmetry as a local
consequence of recursive stabilization.

Time does not reverse. Structures persist.

%% ══════════════════════════════════════════════════════════════════════════
\section{Void Dominance and Topological Redistribution}

\subsection{Expansion Without Metric Inflation}

Barbour correctly recognizes that observational cosmology directly observes
changes in relational structure rather than an externally measurable
expansion through a fixed metric backdrop. RSVP extends this critique
further by questioning whether cosmological evolution should be interpreted
primarily as metric expansion at all.

In the RSVP picture, the universe does not fundamentally grow larger.
Rather, its admissibility structure redistributes.

What conventional cosmology interprets as expansion increasingly appears as
a shift in topological dominance: void regions progressively occupy larger
fractions of the relational accessibility landscape while dense structures
become increasingly filamentary and sparse. Galactic superstructures do not
separate because space itself stretches. Instead, recursively admissible
low-density regions become dominant relative to highly constrained baryonic
configurations. The visible universe appears more diffuse because
admissibility pathways favor smoothing and redistribution over terminal
concentration.

The ontological distinction across the three frameworks is therefore:
\begin{align*}
\text{Standard cosmology:}&\quad \text{Space expands} \\
\text{Barbour:}&\quad \text{Shape complexity increases} \\
\text{RSVP:}&\quad \text{Accessibility structure redistributes}
\end{align*}

The universe is not primarily a metric object. It is a recursively
negotiated admissibility field whose large-scale behavior is governed by
the ongoing redistribution of structural coherence.

\subsection{Filamentary Dissolution}

As void dominance increases, galactic filaments become progressively thinner
and less dynamically influential. Dense matter distributions appear
increasingly wispy relative to the growing dominance of low-density
accessibility regions.

This process is not thermodynamic decay in the conventional sense.
It is a redistribution of admissibility gradients. Highly localized
structures lose their capacity to maintain strong causal differentiation
relative to the increasing smoothness of the surrounding plenum.

The universe therefore evolves toward large-scale topological smoothing
without requiring either total collapse or runaway metric expansion.
This avoids two standard cosmological endpoints simultaneously: the
gravitational singularity of terminal collapse, and the irreversible
dilution implied by metric heat death.

Instead, RSVP predicts progressive relational homogenization through
admissibility redistribution — a smoothing that is neither catastrophic
nor entropic in the classical sense.

\begin{figure}[H]
\centering
\begin{tikzpicture}[scale=0.88]
  % Early dense filamentary structure (left)
  \draw[thick]
    (-5.5,0) -- (-4.5,0.9) -- (-3.6,0.2) -- (-2.8,1.1) -- (-2.0,0);
  \foreach \x/\y in {-5.5/0,-4.5/0.9,-3.6/0.2,-2.8/1.1,-2.0/0} {
    \filldraw (\x,\y) circle (2.5pt);
  }
  \node[font=\small\itshape, above=14pt] at (-3.8,0.9)
        {dense filamentary epoch};

  % Transition arrow
  \draw[-{Latex[length=4mm]}, thick] (-1.4,0.5) -- (0.2,0.5);
  \node[font=\small, above=2pt] at (-0.6,0.5) {$t \to \infty$};

  % Late wispy structure (right)
  \draw[dashed, gray]
    (0.8,0) -- (1.8,0.35) -- (2.9,0.15);
  \foreach \x/\y in {0.8/0,1.8/0.35,2.9/0.15} {
    \filldraw[gray] (\x,\y) circle (1.5pt);
  }

  % Void regions
  \draw[dotted, gray!70] (1.8,1.5) circle (1.25);
  \draw[dotted, gray!70] (2.2,-1.3) circle (1.55);
  \node[font=\small\itshape, gray] at (1.8,3.0) {void dominance};
\end{tikzpicture}
\caption{Topological redistribution in RSVP cosmology. Dense filamentary
structure (left) becomes progressively attenuated as low-density
admissibility regions (dotted) dominate the accessible relational landscape.
The transition is not metric expansion but structural redistribution.}
\end{figure}

%% ══════════════════════════════════════════════════════════════════════════
\section{Poincar\'e Smoothing and the Return of Symmetry}

\subsection{Smoothness Without Maximum Disorder}

One of the deepest differences between RSVP cosmology and both standard
thermodynamics and Barbour's entaxy framework concerns the asymptotic fate
of structure.

Barbour treats increasing complexity as effectively monotonic. The universe
moves away from maximal uniformity toward progressively richer clustering
and differentiation, with no developed principle of reintegration.

RSVP instead predicts that sufficiently long timescales produce large-scale
admissibility smoothing. After Poincar\'{e}-like recurrence intervals,
recursively unstable distinctions are gradually eliminated. Dense structures
dissolve, filaments attenuate, and admissibility gradients flatten. Yet
this process does not produce maximum disorder in the thermodynamic sense.

Instead, the universe approaches a state of maximal relational redundancy.

This distinction is essential. A thermodynamic heat death corresponds to the
exhaustion of usable energy gradients within a statistical ensemble. RSVP
smoothing instead corresponds to the collapse of meaningful
distinguishability itself: variations become operationally identical because
their admissibility consequences converge. The resulting state resembles a
crystal lattice more closely than a thermal gas.

\subsection{The Black Hole Era as Admissibility Purification}

In conventional cosmology, the black hole era represents the exhaustion of
large-scale free-energy gradients. Galaxies disperse, stars decay, and black
holes eventually evaporate through Hawking radiation, leaving the universe
in an increasingly dilute thermal equilibrium.

RSVP interprets this era differently.

The late cosmological regime is not merely the disappearance of matter
structure. It is the progressive elimination of unstable admissible
distinctions. Black holes act not as entropy sinks in the conventional
thermodynamic sense, but as long-duration smoothing operators that remove
recursively unsustainable asymmetries from the admissibility field. They are
better understood as extreme local admissibility concentration regimes
participating in the same redistribution dynamics that governs all structural
evolution within the plenum.

Over cosmological timescales, localized structures lose their capacity to
preserve independent causal differentiation. Filaments attenuate, constraint
gradients flatten, and recursively isolated structures dissolve back into the
surrounding admissibility field.

The universe therefore undergoes a form of admissibility purification.
What disappears is not merely matter concentration, but operationally
meaningful distinction itself. The black hole era does not terminate
cosmological evolution; it acts as a gigantic filtering process through which
distinctions incapable of remaining recursively integrated are progressively
eliminated.

\subsection{Lawfulness Without Structure}

Yet the disappearance of structure does not imply the disappearance of
lawfulness.

In RSVP, laws are not external prescriptions imposed upon matter. They are
the invariant admissibility relations that survive recursive smoothing.
As cosmological evolution eliminates unstable distinctions, progressively
fewer configurations remain operationally differentiable. Matter structures
fade, but admissibility invariants persist.

This produces a crucial inversion of standard thermodynamic logic.

The asymptotic universe is not maximally chaotic. It is maximally redundant.
Remaining configurations increasingly belong to the same admissibility
equivalence classes because their future trajectory consequences converge.
The universe progressively forgets particular structures while preserving
the consistency relations governing admissible transformation itself.

\begin{center}
\itshape
Objects disappear.\\
Persistence survives.
\end{center}

The final cosmological state therefore retains lawfulness even after nearly
all conventional structure has dissolved. Laws persist not because they are
encoded in surviving objects, but because they are the invariant relational
constraints that recursive smoothing cannot eliminate.

In a sense, the universe ends not as matter but as consistency.



In the asymptotic RSVP regime, the universe approaches an effectively
scale-invariant smoothing state in which any remaining variations become
recursively equivalent under admissibility transformations. Distinctions no
longer meaningfully alter future trajectory accessibility.

This produces an unusual inversion of entropy logic.

Conventionally, maximal smoothness is associated with maximal entropy. Within
RSVP, however, a perfectly recursive lattice possesses extremely low effective
entropy precisely because admissible distinctions have collapsed into
equivalence classes. The state is not one of maximal disorder but of maximal
informational gauge redundancy.

The universe does not end in chaos. It ends in recursive redundancy.

This asymptotic state is not a return to initial conditions in any cyclical
cosmological sense, nor is it a reversal through a Janus pivot. The
late-universe lattice is historically generated through recursive smoothing
and admissibility redistribution. Its symmetry is emergent rather than
primordial — the residue of long-duration constraint selection rather than
a recapitulation of origin.

\subsection{Geometry as the Residue of Persistence}

This clarifies the deepest inversion separating RSVP from Barbour's
relational cosmology.

For Barbour, geometry generates persistence: the relational structure of
shape space determines how complexity evolves, and persistence is what
follows from dynamical stability within that geometry.

For RSVP, recursive persistence generates effective geometry: stable
geometric structure is the stabilized shadow left behind by admissibility
relations that survived recursive constraint filtering across cosmological
timescales.

Geometry is therefore not fundamental ontology. It is stabilized persistence
made visible. And smoothness — which for Barbour marks only the improbable
beginning — marks for RSVP both the beginning and the asymptotic limit of
cosmological evolution: two faces of the same admissibility condition,
separated by the entire history of recursive structure formation.

%% ══════════════════════════════════════════════════════════════════════════
\section{Conclusion: Admissibility as the Architecture of Persistence}

The transition from Barbour's relational cosmology to RSVP's admissibility
geometry represents a genuine shift in the philosophy of physics — not merely
an alternative cosmological model, but a reconception of what physics is for.

Barbour dismantles many inherited assumptions: absolute scale, absolute
temporal background, externally imposed reference frames. His relational
reduction is philosophically serious and technically precise. The Janus Point
represents a real contribution to understanding how temporal asymmetry can
emerge without stipulated initial conditions.

RSVP extends this dismantling by questioning the status of physical law
itself:

In standard cosmology, laws generate trajectories; in Barbour,
relational geometry generates temporal structure; in RSVP, recursive
admissibility generates effective lawfulness.

Laws become emergent operators of consistency management rather than
transcendent prescriptions imposed upon matter. The universe increasingly
resembles a self-organizing protocol for preserving recursively stable
structure across scales — not because it "wants" structure, but because
only structure that remains admissible persists long enough to appear
as structure at all.

The Janus Point survives within this framework not as a universal origin,
but as a local admissibility extremum — a region of the landscape where
differentiation pressure has temporarily stabilized near maximum symmetry.
Barbour's shape space survives as a special case of the projection
formalism, with the crucial addition that compression is dynamically
consequential. His geometric records survive as operational constraints,
with the crucial addition that memory is infrastructural rather than
retrospective. His bi-universe does not survive: time does not reverse
because directional continuity is a local consequence of recursive
stabilization, not a global consequence of departure from a privileged pivot.

The universe RSVP describes is neither the thermodynamic machine running
toward equilibrium nor the geometric expansion process trending toward
complexity. It is a continuously negotiated admissibility field that begins
in approximate smoothness, differentiates through lamphron pressure, sustains
structure through recursive admissibility management, redistributes
accessibility as voids dominate and filaments attenuate, and asymptotically
approaches a state of maximal relational redundancy — a crystal-lattice
limit whose symmetry is emergent rather than primordial.

Smoothness is both beginning and limit. Between them lies the entire history
of recursive structure: the conditions for cosmological stability, cognitive
coherence, and temporal experience turn out to be, at the appropriate level
of abstraction, the same admissibility conditions — encountered at different
scales of the same ongoing negotiation.

\bigskip
\begin{center}
  \rule{60mm}{0.3pt}
\end{center}

%% ══════════════════════════════════════════════════════════════════════════
\newpage
\section{Further Directions}

The foregoing argument establishes the philosophical architecture of RSVP
admissibility theory. Several directions suggest themselves for rigorous
extension of the framework into mathematically investigable territory. The
following subsections identify the most consequential of these.

\subsection{Admissibility as a Stability Functional}

One of the unresolved questions in any recursive admissibility framework is
whether admissibility can be formalized as a genuine dynamical functional
rather than treated merely as a descriptive heuristic.

The RSVP framework implicitly assumes that recursively stable structures
occupy regions of trajectory space that minimize destructive inconsistency
while preserving coherent future accessibility. This suggests the existence
of a generalized admissibility functional \(\mathcal{L}[\Afield]\) whose
evolution governs whether a configuration remains dynamically maintainable
across scales.

If such a functional exists, recursive persistence begins to resemble a
Lyapunov-like stability principle:
\[
\frac{d\mathcal{L}}{dt} \leq 0.
\]

The significance of this possibility is substantial. It would imply that
recursive admissibility is not simply a philosophical reinterpretation of
survival, but a mathematically identifiable stability condition constraining
which trajectories remain physically realizable. Establishing whether
\(\mathcal{L}[\Afield]\) admits monotonic evolution — or identifying the
conditions under which it fails to do so — would connect the admissibility
flow directly to gradient-flow systems, renormalization group flows, and
Lyapunov stability theory. In this picture, cosmological persistence,
cognitive coherence, and measurement stability all become manifestations of
the same underlying principle: the minimization of recursively destructive
inconsistency.

\subsection{Projection Blindness and Irrecoverable Compression}

The RSVP projection operator does more than reduce representational
complexity. It actively determines which futures remain operationally
accessible.

Once a trajectory space has been compressed into an admissibility manifold
\(\proj : \Xtraj \to \Madm\), the discarded dimensions are not merely
hidden — they may become operationally unrecoverable. The kernel of the
projection,
\[
\ker(\proj),
\]
therefore represents not mathematical redundancy but lost trajectory
accessibility: regions of configuration space that can no longer participate
in future stabilization because the system's projection grammar has filtered
them out.

This has implications far beyond cosmology. Scientific paradigms, cognitive
models, and measurement systems all generate projection-induced blindness.
Once stabilization occurs within a compressed explanatory manifold, future
reasoning inherits the admissibility structure of that compression. The
system no longer merely describes reality; it operates within the
constraints produced by the projection through which reality became legible.
Formalizing \(\ker(\proj)\) as an information-destroying morphism would give
precise content to notions of measurement collapse, historical
irreversibility, and paradigm lock-in.

\subsection{Coherence Before Reconstruction}

Conventional reconstruction-based architectures assume that coherence is a
property recovered only after sufficient aggregation. Signal-processing
systems gather samples into large buffers, perform global transforms, and
infer structure retrospectively from accumulated data.

The RSVP framework proposes the inverse ordering. Coherence exists prior to
reconstruction whenever local admissibility relations remain recursively
stabilized. Structure does not emerge because a global transform discovers
it; structure persists because local recursive constraints continuously
preserve admissible relations. The ontological priority reverses:
\[
\text{Standard reconstruction:} \quad
\text{aggregation} \;\to\; \text{coherence}
\]
\[
\text{RSVP localism:} \quad
\text{coherence} \;\to\; \text{stable aggregation}
\]

This distinction is fundamental rather than merely computational. In the
RSVP picture, large-scale order is the macroscopic residue of locally
persistent admissibility relations that never lost coherence in the first
place. The actionable consequence is that toy simulations demonstrating
local recurrence stabilization outperforming global reconstruction under
noise and latency constraints would constitute direct experimental evidence
for this ordering — and cellular-automata-scale models could already test
it.

\subsection{Equivalence Collapse and the Late Universe}

The asymptotic smoothing regime described in Section~X should not be
interpreted as the destruction of order, but as the progressive collapse of
operationally meaningful distinction. As recursive smoothing proceeds,
increasingly many trajectories become admissibly equivalent:
\[
d_{\Afield}(\gamma_1,\gamma_2) \to 0.
\]

This does not imply that all microscopic variation disappears. Rather, it
means that remaining distinctions no longer alter future admissibility
structure in physically meaningful ways. The late universe approaches a
quotient-degenerate regime in which multiple configurations belong to the
same operational equivalence class. The crystal-lattice limit is not perfect
classical order; it is recursive smoothing that has eliminated distinctions
producing different admissibility outcomes. The universe asymptotically
forgets individuality while preserving consistency. Making this precise
requires identifying the quotient structure \(\Omega/G_{\Afield}\) where
\(G_{\Afield}\) denotes the group of admissibility-preserving
transformations — a well-posed mathematical question whose answer would
rigorously distinguish RSVP's asymptotic prediction from thermodynamic
heat death.

\subsection{Geometry as an Emergent Compression Layer}

A central implication of the foregoing framework is that geometry may be
secondary to recursive admissibility rather than foundational to it.

In standard physical ontology, geometry provides the primary stage upon
which matter evolves. In Barbour's framework, geometry becomes relational
but retains explanatory priority through the structure of shape space. RSVP
instead suggests that geometry is the stabilized compression layer produced
by recursively persistent admissibility relations. Stable geometries survive
because they minimize destructive inconsistency across scales. Spatial
structure therefore behaves less like an externally existing arena and more
like a dynamically maintained equivalence structure generated through
recursive stabilization.

The strongest eventual formalization of this claim would likely be
categorical rather than metric: admissibility-preserving morphisms,
gluing conditions across local patches, obstruction cancellation, and
local-to-global consistency theorems are already implicit in the framework's
architecture. In that formulation, geometry, cognition, measurement,
lawfulness, and persistence become special cases of recursive constraint
propagation across admissibility-preserving morphisms — and RSVP becomes
most rigorous not as an alternative cosmology but as a generalized
consistency geometry.

%% ══════════════════════════════════════════════════════════════════════════
\newpage
\appendix
\section*{Appendices}
\addcontentsline{toc}{section}{Appendices}

%% ══════════════════════════════════════════════════════════════════════════
\section{Recursive Admissibility Dynamics}

This appendix sketches a minimal mathematical formalization of recursive
admissibility within the RSVP framework. The goal is not to derive a fully
predictive cosmology, but to specify the structural conditions under which
recursive persistence can emerge as a dynamical invariant.

Let the universe be represented by a trajectory space
\[
\Xtraj = \{ \gamma(t) \mid \gamma : \mathbb{R} \to \Omega \},
\]
where \(\Omega\) denotes the total relational configuration space.

The admissibility field
\[
\Afield : \Omega \to \mathbb{R}_{\geq 0}
\]
assigns to each configuration a measure of recursively maintainable
continuity. Configurations with low admissibility fail to preserve stable
constraint propagation across scales, while high-admissibility regions
support persistent structural coherence.

The dynamical evolution of admissibility is written schematically as
\[
\frac{\partial \Afield}{\partial t}
=
\lambda \nabla^2 \Afield
-
\mu \|\nabla \Afield\|^2
+
\nu \mathcal{C}(\Omega),
\]
where \(\lambda\) governs smoothing and reintegration dynamics
(lamphrodyne pressure), \(\mu\) governs local specialization and constraint
concentration (lamphron pressure), and \(\mathcal{C}(\Omega)\) represents
recursive coherence production.

The competition between the smoothing term and the specialization term
determines whether a configuration remains recursively admissible.

A stable admissibility basin satisfies
\[
\frac{\partial \Afield}{\partial t} \approx 0,
\]
while maintaining bounded curvature:
\[
\sup_{\Omega} \|\nabla^2 \Afield\| < \infty.
\]

Recursive persistence therefore corresponds not to static equilibrium but to
metastable admissibility balance.

The projection operator
\[
\proj : \Xtraj \to \Madm
\]
maps the full trajectory space into an admissibility manifold
\[
\Madm \subseteq \Omega
\]
containing only recursively accessible trajectories.

The physically realized universe is therefore not the total trajectory space
itself but the subset that survives recursive admissibility filtering:
\[
\Madm
=
\{ \gamma \in \Xtraj \mid \Afield(\gamma(t)) > \epsilon
\quad \forall t \in I \},
\]
for some admissibility threshold \(\epsilon > 0\).

This formalizes the central RSVP principle:

\begin{quote}
\itshape
Persistence is not the default state of matter.
It is the subset of trajectories capable of recursively preserving
admissibility across scales.
\end{quote}

%% ══════════════════════════════════════════════════════════════════════════
\section{Projection Geometry and Admissibility Compression}

Barbour's shape space construction removes redundant coordinates by
quotienting relational configurations under translation, rotation, and
scaling symmetries. RSVP generalizes this procedure by treating projection
itself as dynamically consequential.

Let \(\Omega\) denote the total relational configuration space and let
\(G\) be a symmetry group acting on \(\Omega\). The quotient
\[
\Omega / G
\]
defines an equivalence structure in which physically redundant distinctions
are identified. Barbour's shape space may therefore be represented as
\[
\mathbb{S}
=
\Omega / \mathrm{Sim}(3),
\]
where \(\mathrm{Sim}(3)\) denotes the similarity group generated by
translations, rotations, and global scalings.

RSVP extends this construction by introducing admissibility-weighted
projection operators. Instead of treating all quotient directions as equally
irrelevant, the RSVP projection selects trajectories according to recursive
persistence criteria:
\[
\proj(\gamma)
=
\arg\max_{\eta \sim \gamma}
\Afield(\eta),
\]
where \(\eta \sim \gamma\) denotes equivalence under
admissibility-preserving transformations.

Projection therefore ceases to be purely representational. It becomes
dynamically selective. This induces an admissibility metric on trajectory
space:
\[
d_{\Afield}(\gamma_1,\gamma_2)
=
\inf_{\phi \in G}
\int
\bigl|
\Afield(\gamma_1(t))
-
\Afield(\phi \cdot \gamma_2(t))
\bigr|
\, dt.
\]

Two trajectories become effectively indistinguishable when their
admissibility difference vanishes:
\[
d_{\Afield}(\gamma_1,\gamma_2) \to 0.
\]

The asymptotic RSVP universe therefore approaches an admissibility quotient
regime in which increasingly many trajectories collapse into the same
equivalence classes. The late universe is not maximally disordered; it is
maximally quotient-redundant. Distinctions disappear because
admissibility-preserving transformations map remaining trajectories into
operational equivalence classes.

%% ══════════════════════════════════════════════════════════════════════════
\section{Asymptotic Smoothing and the Crystal-Lattice Limit}

Let \(\rho(x,t)\) denote the relational differentiation density of the
plenum. Regions of high \(\rho\) correspond to strong localized structure,
constraint concentration, and causal differentiation. The long-duration
smoothing dynamics may be modeled through a nonlinear admissibility flow:
\[
\frac{\partial \rho}{\partial t}
=
D \nabla^2 \rho
-
\alpha \rho^n
+
\beta \mathcal{R}(\rho),
\]
where \(D \nabla^2 \rho\) represents large-scale smoothing,
\(\alpha \rho^n\) represents the decay of recursively unstable
concentrations, and \(\mathcal{R}(\rho)\) represents recursive
stabilization feedback.

During the structure-dominated epoch,
\[
\beta \mathcal{R}(\rho) > \alpha \rho^n,
\]
allowing galaxies, stars, and coherent structures to persist.

During the asymptotic smoothing regime,
\[
D \nabla^2 \rho + \alpha \rho^n \gg \beta \mathcal{R}(\rho),
\]
and recursively unstable distinctions progressively dissolve.

The system evolves toward an admissibility fixed point
\[
\rho(x,t) \to \rho_\infty, \qquad \nabla \rho_\infty \to 0.
\]

Unlike conventional thermal equilibrium, \(\rho_\infty\) does not correspond
to maximum statistical disorder. Instead, the asymptotic state satisfies an
admissibility equivalence condition:
\[
\forall\, x_i, x_j \in \Omega, \qquad \proj(x_i) = \proj(x_j),
\]
meaning that remaining distinctions no longer alter future admissibility
structure.

The late universe therefore approaches a quotient-symmetric state in which
all remaining admissible configurations belong to the same operational
equivalence class:

\begin{quote}
\itshape
The universe does not end in chaos.
It ends in recursive redundancy.
\end{quote}

The asymptotic cosmological state is therefore best interpreted not as heat
death but as admissibility degeneracy: the collapse of meaningful
distinguishability under recursive smoothing.

\begin{figure}[H]
\centering
\begin{tikzpicture}[scale=1.0]
  % Initial distinct trajectories (top row)
  \foreach \x/\lbl in {-4/x_1, -2/x_2, 0/x_3, 2/x_4, 4/x_5} {
    \filldraw (\x,2.4) circle (2.5pt);
    \node[above=3pt, font=\small] at (\x,2.4) {\(\lbl\)};
  }
  % Convergence arrows
  \foreach \x in {-4,-2,0,2,4} {
    \draw[-{Latex[length=2mm]}, gray!80] (\x,2.2) -- (0,0.6);
  }
  % Final equivalence class
  \draw[thick] (0,0) circle (0.9);
  \node[font=\small] at (0,0) {\(\proj(x_i)\)};
  \node[font=\small\itshape, below=14pt] at (0,0)
        {admissibility equivalence class};
\end{tikzpicture}
\caption{Asymptotic recursive smoothing collapses distinguishable
trajectories \(x_1,\ldots,x_5\) into a shared admissibility equivalence
class. The late universe is not maximally disordered but maximally
quotient-redundant: distinctions cease to produce materially different
admissible futures.}
\end{figure}

%% ══════════════════════════════════════════════════════════════════════════
\newpage
\begin{thebibliography}{99}

\bibitem{Barbour2020}
Julian Barbour,
\textit{The Janus Point: A New Theory of Time},
Basic Books, 2020.

\bibitem{BarbourBertotti1982}
Julian Barbour and Bruno Bertotti,
``Mach's Principle and the Structure of Dynamical Theories,''
\textit{Proceedings of the Royal Society A},
vol.~382, no.~1783, pp.~295--306, 1982.

\bibitem{Barbour2003}
Julian Barbour,
``Scale-Invariant Gravity: Particle Dynamics,''
\textit{Classical and Quantum Gravity},
vol.~20, pp.~1543--1570, 2003.

\bibitem{BarbourKoslowskiMercati2014}
Julian Barbour, Tim Koslowski, and Flavio Mercati,
``Identification of a Gravitational Arrow of Time,''
\textit{Physical Review Letters},
vol.~113, no.~18, 181101, 2014.

\bibitem{Mercati2018}
Flavio Mercati,
\textit{Shape Dynamics: Relativity and Relationalism},
Oxford University Press, 2018.

\bibitem{Mach1893}
Ernst Mach,
\textit{The Science of Mechanics},
Open Court Publishing, 1893.

\bibitem{Boltzmann1896}
Ludwig Boltzmann,
\textit{Lectures on Gas Theory},
University of California Press, 1896.

\bibitem{Penrose1979}
Roger Penrose,
``Singularities and Time-Asymmetry,''
in \textit{General Relativity: An Einstein Centenary Survey},
ed.\ S.~W. Hawking and W.~Israel,
Cambridge University Press, 1979.

\bibitem{Penrose2004}
Roger Penrose,
\textit{The Road to Reality},
Jonathan Cape, 2004.

\bibitem{Hawking1975}
Stephen W. Hawking,
``Particle Creation by Black Holes,''
\textit{Communications in Mathematical Physics},
vol.~43, pp.~199--220, 1975.

\bibitem{Bekenstein1973}
Jacob D. Bekenstein,
``Black Holes and Entropy,''
\textit{Physical Review D},
vol.~7, no.~8, pp.~2333--2346, 1973.

\bibitem{Carroll2010}
Sean Carroll,
\textit{From Eternity to Here},
Dutton, 2010.

\bibitem{Prigogine1984}
Ilya Prigogine and Isabelle Stengers,
\textit{Order Out of Chaos},
Bantam Books, 1984.

\bibitem{Kauffman1993}
Stuart Kauffman,
\textit{The Origins of Order},
Oxford University Press, 1993.

\bibitem{Haken1983}
Hermann Haken,
\textit{Synergetics: An Introduction},
Springer, 1983.

\bibitem{Friston2010}
Karl Friston,
``The Free-Energy Principle: A Unified Brain Theory?''
\textit{Nature Reviews Neuroscience},
vol.~11, pp.~127--138, 2010.

\bibitem{FristonKilnerHarrison2006}
Karl Friston, James Kilner, and Lee Harrison,
``A Free Energy Principle for the Brain,''
\textit{Journal of Physiology-Paris},
vol.~100, pp.~70--87, 2006.

\bibitem{Shannon1948}
Claude E. Shannon,
``A Mathematical Theory of Communication,''
\textit{Bell System Technical Journal},
vol.~27, pp.~379--423 and 623--656, 1948.

\bibitem{Wheeler1990}
John Archibald Wheeler,
``Information, Physics, Quantum: The Search for Links,''
in \textit{Complexity, Entropy, and the Physics of Information},
ed.\ W.~Zurek, Addison-Wesley, 1990.

\bibitem{Landauer1961}
Rolf Landauer,
``Irreversibility and Heat Generation in the Computing Process,''
\textit{IBM Journal of Research and Development},
vol.~5, no.~3, pp.~183--191, 1961.

\bibitem{Smolin2001}
Lee Smolin,
\textit{Three Roads to Quantum Gravity},
Basic Books, 2001.

\bibitem{Rovelli2004}
Carlo Rovelli,
\textit{Quantum Gravity},
Cambridge University Press, 2004.

\bibitem{Rovelli2018}
Carlo Rovelli,
\textit{The Order of Time},
Riverhead Books, 2018.

\bibitem{GellMannHartle1994}
Murray Gell-Mann and James B. Hartle,
``Time Symmetry and Asymmetry in Quantum Mechanics and Quantum Cosmology,''
in \textit{Physical Origins of Time Asymmetry},
ed.\ J.~J. Halliwell et al.,
Cambridge University Press, 1994.

\bibitem{Zurek2003}
Wojciech H. Zurek,
``Decoherence, Einselection, and the Quantum Origins of the Classical,''
\textit{Reviews of Modern Physics},
vol.~75, pp.~715--775, 2003.

\bibitem{Lloyd2006}
Seth Lloyd,
\textit{Programming the Universe},
Knopf, 2006.

\bibitem{Deutsch1997}
David Deutsch,
\textit{The Fabric of Reality},
Penguin Books, 1997.

\bibitem{Whitehead1929}
Alfred North Whitehead,
\textit{Process and Reality},
Macmillan, 1929.

\bibitem{DeLanda2002}
Manuel DeLanda,
\textit{Intensive Science and Virtual Philosophy},
Continuum, 2002.

\bibitem{Anderson1972}
Philip W. Anderson,
``More Is Different,''
\textit{Science},
vol.~177, no.~4047, pp.~393--396, 1972.

\end{thebibliography}

\end{document}

